\section{Conclus\~ao e Trabalhos Futuros}
\label{sec:conclusao}

A pergunta deste trabalho foi se as falas anteriores de uma pessoa, condensadas num perfil, ajudam
um LLM a reconhecer a opinião que a imprensa lhe atribuiu numa audiência posterior. Para isso, o
pipeline liga as opiniões publicadas sobre audiências públicas à fala de quem as emitiu e condensa a
fala de cada ator em um perfil escrito somente a partir dela, usado para simular como a pessoa se
posicionaria.

A UDV localiza evidência candidata para 2.105 das 2.203 opiniões, mas, num controle com as 183
opiniões de treino que contêm citação confiável, a sentença de maior cosseno coincide com a passagem
citada em 61,2\% dos casos. Na validação humana, a precisão estrita da citação literal foi de 0,714
e a tolerante do nível \texttt{semantic\_match\_high}, de 0,738, e nenhum dos dois critérios
registrados foi atingido. Na simulação, o perfil aumentou o acerto da múltipla escolha em 0,139
no Gemma 4 31B e em 0,119 no Qwen3.8 27B, com intervalos de 95\% acima de zero, e os trechos
recuperados não tiveram efeito detectado.

Ficam em aberto uma condição só com os trechos recuperados, a leitura dos perfis de um modelo por
outro e a comparação entre prompts de perfil diferentes. A resposta simulada é uma estimativa condicionada às falas públicas da
pessoa e não equivale a uma declaração dela. Por isso, nunca deve ser apresentada como citação e
não serve para inferir intenção política.
